\section{Билет 1. Колебательное движение. Виды колебаний. Вертикальный пружинный маятник. Уравнение гармонических колебаний (с выводом).}

\begin{definition}
\textbf{Колебаниями} (колебательными движениями) называются движения или процессы, отличающиеся той или иной степенью повторяемости во времени. 
Физическая природа колебаний может быть различной: механические (маятник часов, груз на пружине, колебания молекул), электромагнитные (ток в цепи, радиоволны), звуковые и т.д. Несмотря на различие природы, все они описываются одинаковыми математическими моделями.
\end{definition}

\begin{definition}
В зависимости от характера воздействия на систему различают:
\begin{enumerate}
    \item \textbf{Свободные (собственные) колебания} --- колебания, происходящие в системе, предоставленной самой себе, при отсутствии внешних воздействий после того, как она была выведена из положения устойчивого равновесия или ей был сообщён начальный импульс.
    \item \textbf{Вынужденные колебания} --- колебания, протекающие под действием внешней периодически изменяющейся силы.
\end{enumerate}
\end{definition}

\begin{definition}
\textbf{Гармоническими колебаниями} называются колебания, при которых колеблющаяся величина изменяется со временем по закону синуса или косинуса. Это простейший и фундаментальный вид колебательного движения.
\end{definition}

\begin{theorem}[Уравнение движения вертикального пружинного маятника]
Рассмотрим систему: тело массы $m$, подвешенное на вертикальной пружине жёсткостью $k$. 
\begin{figure}[htbp]
    \centering
    \begin{tikzpicture}[
        scale=1.2,
        spring/.style={decorate, decoration={zigzag, amplitude=4pt, segment length=6pt}},
        mass/.style={fill=blue!20, draw=blue!70, thick},
        mass_eq/.style={fill=gray!15, draw=gray!50, thick, dashed},
        force/.style={->, >=Stealth, thick},
        vec_red/.style={force, red!80!black},
        vec_blue/.style={force, blue!80!black},
        vec_gray/.style={force, gray!70}
    ]
    \def\ydisp{-1.2} % Координата центра смещённого груза
    
    % Потолок
    \fill[gray!30] (-1.5, 3) rectangle (1.5, 3.3);
    \draw[thick] (-1.5, 3) -- (1.5, 3);
    
    % Ось Ox (направлена вниз, начало в равновесии)
    \draw[->, thick] (-2.5, 3) -- (-2.5, -2.5) node[below] {$Ox$};
    \node[left] at (-2.5, 1) {$x=0$};
    \node[left, font=\small] at (-2.5, -1.2) {$x$};
    
    % Положение равновесия (пунктир)
    \draw[spring, gray!50, dashed, thick] (0, 3) -- (0, 1.5);
    \draw[mass_eq] (-0.6, 0.5) rectangle (0.6, 1.5);
    \node[gray!50] at (0, 1) {$m$};
    
    % Силы в равновесии (уравновешены)
    \draw[vec_gray] (0, 1) -- (0, 2.3) node[midway, right, gray] {\small $k\Delta l$};
    \draw[vec_gray] (0, 1) -- (0, -0.3) node[midway, right, gray] {\small $mg$};
    
    % Смещённое положение (сплошной цвет)
    \draw[spring, blue!70!black, thick] (0, 3) -- (0, \ydisp+0.5);
    \draw[mass] (-0.6, \ydisp-0.5) rectangle (0.6, \ydisp+0.5);
    \node at (0, \ydisp) {$m$};
    
    % Статическое удлинение
    \draw[<->, gray, thin] (0.8, 3) -- (0.8, 1) node[midway, right, gray, font=\small] {$\Delta l$};
    
    % Смещение x
    \draw[<->, thick, purple] (1.2, 1) -- (1.2, \ydisp) node[midway, right, fill=white, font=\bfseries, purple] {$x$};
    \node[right, align=left, font=\small] at (1.5, -0.2) {смещение\\из равновесия};
    
    % Силы в смещённом положении
    % Сила упругости (вверх, длиннее mg, т.к. пружина растянута сильнее)
    \draw[vec_red] (0, \ydisp+0.3) -- (0, \ydisp+1.6) node[midway, right] {$F_{упр}$};
    % Сила тяжести (вниз)
    \draw[vec_red] (0, \ydisp-0.3) -- (0, \ydisp-1.4) node[midway, right] {$mg$};
    % Результирующая возвращающая сила (вверх, слева)
    \draw[vec_blue] (-1.2, \ydisp) -- (-1.2, \ydisp+0.9) node[midway, left, font=\bfseries, blue] {$F=-kx$};
    
    \end{tikzpicture}
    \caption{Вертикальный пружинный маятник. Серым пунктиром показано положение статического равновесия ($x=0$), синим --- текущее смещённое положение. Ось $Ox$ направлена вниз.}
    \label{fig:vertical_spring}
\end{figure}


Направим ось $Ox$ вниз, начало координат поместим в положение статического равновесия тела. В этом положении сила тяжести уравновешивается силой упругости:
\begin{equation}
    mg = k \Delta l, \label{eq:equilibrium}
\end{equation}
где $\Delta l$ --- статическое удлинение пружины.

Если сместить тело из положения равновесия на величину $x$, то результирующая сила, действующая на тело, с учётом закона Гука будет равна:
\begin{equation}
    F = mg - k(\Delta l + x).
\end{equation}
Подставляя \eqref{eq:equilibrium}, получаем:
\begin{equation}
    F = -kx. \label{eq:restoring}
\end{equation}
Знак «минус» указывает на то, что сила всегда направлена к положению равновесия (возвращающая сила).

Запишем второй закон Ньютона в проекции на ось $x$:
\begin{equation}
    m \frac{d^2x}{dt^2} = -kx.
\end{equation}
Перенеся все слагаемые в одну сторону и разделив на $m$, получаем линейное однородное дифференциальное уравнение второго порядка:
\begin{equation}
    \ddot{x} + \frac{k}{m}x = 0.
\end{equation}
Обозначив $\omega_0^2 = \frac{k}{m}$, приходим к каноническому виду уравнения гармонических колебаний:
\begin{equation}
    \ddot{x} + \omega_0^2 x = 0. \label{eq:harmonic_eq}
\end{equation}
\end{theorem}

\begin{remark}
Величина $\omega_0 = \sqrt{k/m}$ называется \textbf{циклической (круговой) частотой} собственных колебаний. Силы вида $F = -kx$, имеющие неупругую природу (например, гравитационные или электростатические в малых приближениях), называют \textbf{квазиупругими}. Уравнение \eqref{eq:harmonic_eq} универсально для всех систем с квазиупругой возвращающей силой.
\end{remark}
